% Twelfth Philippic --- headnote
% philippic-12, March 43 BC, Rome

\textit{Delivered in the Senate early in March 43 BC, the Twelfth
Philippic is a speech of retraction. The war at Mutina was still
undecided: Decimus Brutus lay besieged, the consuls Hirtius and
Pansa were taking the field, and the young Caesar (Octavian) was in
arms beside them. Into this the two consulars Lucius Calpurnius Piso
and Quintus Fufius Calenus --- both intimates of Antony --- had
revived the hope of an ``honourable peace,'' and the consul Pansa
had lent it his weight; a second embassy to Antony was resolved upon,
and Cicero, against his own deeper judgment, had let himself be named
to it. (An earlier embassy of three --- one of whom, Servius
Sulpicius Rufus, had died on the road --- had already failed.) The
present speech is Cicero's withdrawal, both of his assent to the
embassy and of his own person from it.}

\textit{The argument moves on two fronts. First, that the embassy is
folly: Cicero confesses he was deceived, dazzled (he says) by his
care for Decimus Brutus, and caught by Fufius Calenus's question
whether Antony would not be heard even if he yielded; second
thoughts, he answers, are the wiser, and ``the best harbour for a man
who repents is a change of plan.'' A peace mission can do the
commonwealth no good and much harm: it slackens the recovered
longing of the towns and colonies and of all Italy, cuts the sinews
of the legions --- the heaven-sent Martian and the Fourth --- and
binds the hands of the very commanders, Hirtius and the young Caesar,
whose letters in his hand promise victory. And in any case no peace
is possible: every concession would unmake the Senate's own decrees
--- the forged minutes, the laws carried by violence, the seven
hundred million sesterces embezzled, the offices and kingdoms sold
--- and would hand the war's victory to Antony for nothing. To make
peace with Antony's brigands, he says, ``will not be peace, but a
compact of slavery.''}

\textit{The second front is personal and gives the speech its
peculiar colour. Cicero --- whom Antony's harangues had marked out as
his bitterest private enemy, and whose confiscated goods Antony had
been parcelling out to his creatures --- argues that he, of all men,
was the last who should have been chosen to procure such a peace: he
who first put on the soldier's cloak, who always named Antony enemy
and not adversary, war and not tumult. There follows a remarkable
survey of the three roads to Mutina, each (he shows) commanded by one
of Antony's partisans, and a refusal to entrust himself to them ---
he who at the Terminalia had not dared even to cross into the suburbs.
He recalls the famous parleys of the Social War and of Sulla's day
(Pompeius Strabo with Vettius Scato, ``a guest by goodwill, a foe by
necessity''; Sulla with Scipio), as cases where a conference was at
least safe, and judges that no such safety could be had with Antony.
His conclusion is conditional and characteristically self-effacing:
let his life be kept for the commonwealth; he will measure the whole
question not by his own peril but by the public good, and go only if
it can be done safely. In the event the second embassy was abandoned,
and the issue was left, as Cicero wished, to the sword at Mutina.}
